
\section{连接轨道模空间的定向}
\label{sec:8-7}

在第~\ref{subsec:8-1-3} 小节中，我们已经描述了连接轨道模空间上的
定向。\liftIndex{I304}本节将给带标记点伪全纯条带的模空间及其与
$L^{(i)}$ 上光滑奇异单形的纤维积赋予定向。它们用于定义滤过
$A_\infty$ 双模结构。

设 $(M,\omega)$ 中的一对拉格朗日子流形 $(L^{(0)},L^{(1)})$ 相对
spin，并且彼此横截相交。（见定义~\ref{def:8-1-2} 中关于这对子流形的
相对 spin 结构之定义。）对每个 $p\in L^{(0)}\cap L^{(1)}$，选定一条
定向拉格朗日线性子空间道路 $\lambda_p$，并选定丛
$\widetilde\lambda_p$ 的一个平凡化（见第~\ref{subsec:8-1-3} 小节）。
从现在起，在模空间的记号中省略这些资料。

现取
$[\ell_p,w],[\ell_q,w']\in\operatorname{Cr}(L^{(1)},L^{(0)})$
（集合 $\operatorname{Cr}(L^{(1)},L^{(0)})$ 的定义见第 2.2.2 小节）。
记
\[
 \widetilde{\mathcal M}^{\mathrm{reg}}
 ([\ell_q,w'],[\ell_p,w])
 =\widetilde{\mathcal M}^{\mathrm{reg}}
 (L^{(1)},L^{(0)};[\ell_q,w'],[\ell_p,w])
\]
为从无穷圆柱 $\mathbb R\times[0,1]\subset\mathbb C$ 到 $M$ 的伪全纯
映射 $u$ 所成的空间，其中
\[
 \begin{gathered}
 u(-\infty,t)=p,\qquad u(+\infty,t)=q,\qquad
 u(\mathbb R\times\{i\})\subset L^{(i)}\quad(i=0,1),\\
 w\#u\sim w'.
 \end{gathered}
\]
注意，这与第 3.7 节使用的记号不同；在那里，上述模空间采用相反的输入
次序，记为
\[
 \widetilde{\mathcal M}^{\mathrm{reg}}
 ([\ell_p,w],[\ell_q,w']).
\]
回忆，我们只在处理定向时采用这里的不同记号；见第
\ref{subsec:8-1-3} 小节开头。（为简便起见，我们常在
$\widetilde{\mathcal M}(L^{(1)},L^{(0)};
[\ell_q,w'],[\ell_p,w])$ 等记号中省略 $L^{(0)}$ 和 $L^{(1)}$。）

从现在起，把无穷圆柱 $\mathbb R\times[0,1]$ 与
$D^2\setminus\{\pm1\}\subset\mathbb C$ 识别，使端
$\{+\infty\}\times[0,1]$ 与 $\{-\infty\}\times[0,1]$ 分别对应
$z_0=1$ 与 $z_1=-1$。于是，由第 8.1 节定理~\ref{thm:8-1-14}，
$(L^{(0)},L^{(1)})$ 的相对 spin 结构确定
$\widetilde{\mathcal M}^{\mathrm{reg}}
([\ell_q,w'],[\ell_p,w])$ 上的定向。

置
\[
 \partial_1D^2=\{z\in\partial D^2\mid\operatorname{Im}z>0\},
 \qquad
 \partial_0D^2=\{z\in\partial D^2\mid\operatorname{Im}z<0\}.
\]
以
$\widetilde{\mathcal M}^{\mathrm{reg}}_{\ell,m}
([\ell_q,w'],[\ell_p,w])$
表示如下连接轨道
\[
 u\colon(D^2\setminus\{\pm1\};\partial_1D^2,\partial_0D^2)
 \longrightarrow(M;L^{(1)},L^{(0)})
\]
所成的空间：$\partial_1D^2$ 上带有 $\ell$ 个标记点，
$\partial_0D^2$ 上带有 $m$ 个标记点，并且 $u(-1)=p$、$u(1)=q$。
这些标记点彼此不同。第 $2,3,\ldots,\ell+1$ 个标记点
$z_2,\ldots,z_{\ell+1}$ 位于 $\partial_1D^2$ 上，第
$\ell+2,\ldots,\ell+m+1$ 个标记点
$z_{\ell+2},\ldots,z_{\ell+m+1}$ 位于 $\partial_0D^2$
上\translatornote{源书把后一组标记点写成位于 $L^{(0)}$ 上；但这些
$z_i$ 是定义域边界上的点，而映射 $u$ 才把 $\partial_0D^2$ 映入
$L^{(0)}$，故此处恢复为 $\partial_0D^2$。}。
若置 $z_0=1,z_1=-1$，则标记点的次序
\[
 z_0,z_2,\ldots,z_{\ell+1},z_1,z_{\ell+2},\ldots,z_{\ell+m+1}
\]
与 $\partial D^2$ 的逆时针定向相容（见注~\ref{rem:3-7-23}(2)
和图~\ref{fig:3-7-2}）。

空间
$\widetilde{\mathcal M}^{\mathrm{reg}}_{\ell,m}
([\ell_q,w'],[\ell_p,w])$
是
\[
 \widetilde{\mathcal M}^{\mathrm{reg}}
 ([\ell_q,w'],[\ell_p,w])\times(\partial D^2)^{\ell+m}
\]
的一个开子集。因此，在上述约定下，它的定向由
\[
 \widetilde{\mathcal M}^{\mathrm{reg}}_{\ell,m}
 ([\ell_q,w'],[\ell_p,w])
 \subset
 \widetilde{\mathcal M}^{\mathrm{reg}}
 ([\ell_q,w'],[\ell_p,w])
 \times
 (\partial D_2^2\times\cdots\times\partial D_{\ell+m+1}^2)
\]
给出；这里 $\partial D_i^2$ 是标记点 $z_i$
（$i=2,\ldots,\ell+m+1$）的参数空间。

再以 $\mathcal M_{\ell,m}([\ell_q,w'],[\ell_p,w])$ 表示模空间
$\widetilde{\mathcal M}_{\ell,m}([\ell_q,w'],[\ell_p,w])$ 关于固定两个
标记点 $z_0=+1,z_1=-1$ 的 $D^2$ 双全纯自同构群 $\mathbb R$ 所得的
商空间（见约定~\ref{conv:8-3-1}）。其定向定义为第 8.2 节所定义的
商定向。像第 8.3 节一样，我们简写为
\[
 \widetilde{\mathcal M}^{\mathrm{reg}}_{\ell,m}
 ([\ell_q,w'],[\ell_p,w])
 =\mathcal M^{\mathrm{reg}}_{\ell,m}
 ([\ell_q,w'],[\ell_p,w])\times\mathbb R.
\]

令
\[
 P_1^{(1)},\ldots,P_\ell^{(1)}\in C(L^{(1)};\mathbb Q),
 \qquad
 P_1^{(0)},\ldots,P_m^{(0)}\in C(L^{(0)};\mathbb Q)
\]
为 $L^{(i)}$（$i=0,1$）上的光滑奇异单形。

\begin{definition}{8.7.1}\label{def:8-7-1}
定义
\[
\begin{aligned}
 &\mathcal M^{\mathrm{reg}}_{\ell,m}
 \bigl(L^{(1)},L^{(0)};[\ell_q,w'],[\ell_p,w]:
 P_1^{(1)},\ldots,P_\ell^{(1)};
 P_1^{(0)},\ldots,P_m^{(0)}\bigr)\\
 &\quad=
 \mathcal M^{\mathrm{reg}}_{\ell,m}
 \bigl([\ell_q,w'],[\ell_p,w]:
 P_1^{(1)},\ldots,P_\ell^{(1)};
 P_1^{(0)},\ldots,P_m^{(0)}\bigr)\\
 &\quad:=(-1)^\epsilon
 \mathcal M^{\mathrm{reg}}_{\ell,m}
 ([\ell_q,w'],[\ell_p,w])_{
 (\operatorname{ev}_2,\ldots,\operatorname{ev}_{\ell+1},
  \operatorname{ev}_{\ell+2},\ldots,
  \operatorname{ev}_{\ell+m+1})}\\[-2pt]
 &\hspace{13em}{}\times
 \left(\prod_{k=1}^{\ell}P_k^{(1)}
 \times\prod_{k=1}^{m}P_k^{(0)}\right),
\end{aligned}
\]
其中 $\epsilon$ 如下定义。令 $\mu([\ell_p,w])$ 为定义
\ref{def:2-2-12} 中的马斯洛夫--莫尔斯指标。（由于
$\mu([\ell_p,w];\lambda^0)$ 的奇偶性不依赖于 $\lambda^0$，故在记号中
省去 $\lambda^0$；它也不依赖于 $w$。）当 $m=0$ 时，置
\[
\begin{aligned}
 \epsilon={}&(n+1)\sum_{k=1}^{\ell-1}\sum_{j=1}^{k}
 \deg P_j^{(1)}
 +\ell(n+1)\mu([\ell_p,w])+\mu([\ell_q,w'])\\
 &{}+(\mu([\ell_p,w])+1)
 \sum_{k=1}^{\ell}(\deg P_k^{(1)}+1).
\end{aligned}
\]
当 $m\geq1$ 时，置%
\translatornote{源书这里写“$m>1$”，从而没有定义 $m=1$ 的符号。
定义~\ref{def:8-4-1} 的统一标记点规则对任意非空的下边界输入组均适用，
而紧随的公式在 $m=1$ 时各空和也有定义；故恢复为 $m\geq1$。}
\[
\begin{aligned}
 \epsilon={}&(n+1)\Biggl(
 \sum_{k=1}^{\ell}\sum_{j=1}^{k}\deg P_j^{(1)}\\[-2pt]
 &\qquad{}+\sum_{k=1}^{m-1}\left(
   \sum_{h=1}^{\ell}\deg P_h^{(1)}
   +\sum_{j'=1}^{k}\deg P_{j'}^{(0)}\right)\Biggr)\\
 &{}+\mu([\ell_q,w'])
 +(\mu([\ell_p,w])+1)
 \sum_{k=1}^{\ell}(\deg P_k^{(1)}+1)\\
 &{}+(\ell+m)(n+1)\mu([\ell_p,w]).
\end{aligned}
\]
\end{definition}

\begin{remark}{8.7.2}\label{rem:8-7-2}
\begin{enumerate}[label=\textup{(\arabic*)},leftmargin=2.8em]
\item 把连接轨道空间看作以单位圆盘为定义域的伪全纯映射空间时，
这个符号正是由定义~\ref{def:8-4-1} 的规则所确定的符号。注意，我们处理
头两个标记点 $z_0,z_1$ 的方式不同于处理其余标记点的方式。这是因为，
在定义空间 $\widetilde{\mathcal M}([\ell_q,w'],[\ell_p,w])$ 的定向时，
必须如此处理：我们在这两个标记点处胶合半圆盘，胶合后这两个标记点便消失。

\item $L^{(0)}$ 与 $L^{(1)}$ 干净相交的情形将在第 8.8 节讨论。
\end{enumerate}
\end{remark}

空间
\[
 \mathcal M_{\ell,m}
 \bigl([\ell_q,w'],[\ell_p,w]:
 P_1^{(1)},\ldots,P_\ell^{(1)};
 P_1^{(0)},\ldots,P_m^{(0)}\bigr)
\]
是
\[
 \mathcal M^{\mathrm{reg}}_{\ell,m}
 \bigl([\ell_q,w'],[\ell_p,w]:
 P_1^{(1)},\ldots,P_\ell^{(1)};
 P_1^{(0)},\ldots,P_m^{(0)}\bigr)
\]
的稳定映射紧化，并且是带角的仓西结构空间。注意，仓西结构的定向丛
自然地延拓到稳定映射紧化。因此，
\[
 \mathcal M_{\ell,m}
 \bigl([\ell_q,w'],[\ell_p,w]:
 P_1^{(1)},\ldots,P_\ell^{(1)};
 P_1^{(0)},\ldots,P_m^{(0)}\bigr)
\]
也典范地带有定向。利用这些有向零维模空间，像定义
\ref{def:3-7-41} 那样，以
\[
\begin{aligned}
 &\left\langle
 \mathfrak n_{\ell,m}
 \bigl(P_1^{(1)}\otimes\cdots\otimes P_\ell^{(1)}
 \otimes[\ell_p,w]\otimes
 P_1^{(0)}\otimes\cdots\otimes P_m^{(0)}\bigr),
 [\ell_q,w']\right\rangle\\
 &\qquad:=\#\mathcal M_{\ell,m}
 \bigl([\ell_q,w'],[\ell_p,w]:
 P_1^{(1)},\ldots,P_\ell^{(1)};
 P_1^{(0)},\ldots,P_m^{(0)}\bigr)
\end{aligned}
\]
定义 $\mathfrak n_{\ell,m}$。于是可证明下列命题。

\begin{proposition}{8.7.3}\label{prop:8-7-3}
算子 $\mathfrak n_{\ell,m}$ 满足 $A_\infty$ 双模公式。
\end{proposition}

\begin{proof}
模空间
\begin{equation}
 \mathcal M_{\ell,m}
 \bigl([\ell_q,w'],[\ell_p,w]:
 P_1^{(1)},\ldots,P_\ell^{(1)};
 P_1^{(0)},\ldots,P_m^{(0)}\bigr)
 \tag{8.7.4}\label{eq:8-7-4}
\end{equation}
有三类端。全纯圆盘冒泡所对应的端，其处理方式与命题
\ref{prop:8-5-1}(1) 的证明类似。（这些贡献同 $L^{(i)}$（$i=0,1$）
上的 $A_\infty$ 代数运算相消。）可以看出，其余边界贡献两两相消。

我们有
\begin{equation}
\begin{aligned}
 &(-1)^{\epsilon_1}
 \mathcal M_{\ell,m}
 \bigl([\ell_q,w'],[\ell_p,w]:
 P_1^{(1)},\ldots,\partial P_j^{(1)},\ldots,P_\ell^{(1)};
 P_1^{(0)},\ldots,P_m^{(0)}\bigr)\\
 &\qquad\subset
 \partial\mathcal M_{\ell,m}
 \bigl([\ell_q,w'],[\ell_p,w]:
 P_1^{(1)},\ldots,P_\ell^{(1)};
 P_1^{(0)},\ldots,P_m^{(0)}\bigr),
\end{aligned}
\tag{8.7.5}\label{eq:8-7-5}
\end{equation}
其中
\[
 \epsilon_1=\mu([\ell_q,w'])+n+1
 +\sum_{i=1}^{j-1}(\deg P_i^{(1)}+1).
\]
此外，
\begin{equation}
\begin{aligned}
 &(-1)^{\epsilon_2}
 \mathcal M_{\ell,m}
 \bigl([\ell_q,w'],[\ell_p,w]:
 P_1^{(1)},\ldots,P_\ell^{(1)};
 P_1^{(0)},\ldots,\partial P_j^{(0)},\ldots,P_m^{(0)}\bigr)\\
 &\qquad\subset
 \partial\mathcal M_{\ell,m}
 \bigl([\ell_q,w'],[\ell_p,w]:
 P_1^{(1)},\ldots,P_\ell^{(1)};
 P_1^{(0)},\ldots,P_m^{(0)}\bigr),
\end{aligned}
\tag{8.7.6}\label{eq:8-7-6}
\end{equation}
其中
\[
 \epsilon_2=\mu([\ell_q,w'])+n+1
 +\sum_{i=1}^{\ell}(\deg P_i^{(1)}+1)
 +(\mu([\ell_p,w])+1)
 +\sum_{i=1}^{j-1}(\deg P_i^{(0)}+1).
\]
式~\eqref{eq:8-7-5} 和
\eqref{eq:8-7-6} 的证明见下一节的
(8.8.12.2)、(8.8.12.3) 及定理 8.8.10(2) 的证明；在那里，把这些结果
同 (8.8.8) 及定义 8.8.11 结合起来，将证明更一般的陈述。

模空间~\eqref{eq:8-7-4} 的另一些端对应于一条连接轨道分裂为两条连接
轨道之和。对第~\ref{subsec:8-1-3} 小节和约定~\ref{conv:8-6-1}
所给的定向，可以证明胶合映射
\[
 \widetilde{\mathcal M}([\ell_q,w'],[\ell_r,w''])
 \times
 \widetilde{\mathcal M}([\ell_r,w''],[\ell_p,w])
 \longrightarrow
 \widetilde{\mathcal M}([\ell_q,w'],[\ell_p,w])
\]
保定向。%
\translatornote{源书把胶合映射的目标写成
$\widetilde{\mathcal M}([\ell_q,w'],[\ell_p,w''])$；但第二因子的负端资料
是 $[\ell_p,w]$，胶合后仍应保留它，且紧随的取商边界式也以
$[\ell_p,w]$ 为负端，故改正如上。}
注意，这里使用的是第~\ref{subsec:8-1-3} 小节中的定向；定义
\ref{def:8-7-1} 对这些定向作了修正。

现在考虑由 $\tau$-变量平移定义的 $\mathbb R$-作用。依第 8.2 节的约定，
有下列引理~\ref{lem:8-3-5} 的类似结果：
\[
 (-1)^{\epsilon_3}
 \mathcal M([\ell_q,w'],[\ell_r,w''])
 \times\mathcal M([\ell_r,w''],[\ell_p,w])
 \subset
 \partial\mathcal M([\ell_q,w'],[\ell_p,w]),
\]
其中
\[
 \epsilon_3=\mu([\ell_q,w'])-\mu([\ell_r,w''])-1.
\]
在定理 8.8.10(3) 的证明中将会见到一个更一般的陈述。利用上述事实，
得到
\[
\begin{aligned}
 &(-1)^{\epsilon_4}
 \mathcal M_{\ell_1,m_1}
 \bigl([\ell_q,w'],[\ell_r,w'']:
 P_1^{(1)},\ldots,P_{\ell_1}^{(1)};
 P_1^{(0)},\ldots,P_{m_1}^{(0)}\bigr)\\
 &\quad{}\times
 \mathcal M_{\ell_2,m_2}
 \bigl([\ell_r,w''],[\ell_p,w]:
 P_1^{(1)'},\ldots,P_{\ell_2}^{(1)'};
 P_1^{(0)'},\ldots,P_{m_2}^{(0)'}\bigr)\\
 &\quad\subset
 \partial\mathcal M_{\ell,m}
 \bigl([\ell_q,w'],[\ell_p,w]:
 P_1^{(1)},\ldots,P_{\ell_1}^{(1)},
 P_1^{(1)'},\ldots,P_{\ell_2}^{(1)'};\\[-2pt]
 &\hspace{13em}
 P_1^{(0)'},\ldots,P_{m_2}^{(0)'},
 P_1^{(0)},\ldots,P_{m_1}^{(0)}\bigr),
\end{aligned}
\]
其中
\[
 \epsilon_4=\mu([\ell_q,w'])+1
 +\sum_{i=1}^{\ell_1}(\deg P_i^{(1)}+1);
\]
参见 (8.8.12.4)。其余论证现在已是标准的。
\end{proof}
